\chapter{Orbitales de fragments et molécules polyatomiques}\label{ch:b2:fragment-orbitals}

Le schéma de Lewis de l'eau montre deux liaisons O--H équivalentes et deux
doublets non liants équivalents. Pourtant, lorsque l'eau est ionisée par un rayonnement ultraviolet, les
électrons sortent avec quatre énergies différentes, et non deux~; le méthane, avec quatre
liaisons C--H équivalentes, en montre deux. Dans une molécule, les électrons ne se logent pas dans des liaisons
entre paires d'atomes~: ils occupent des orbitales étendues à toute la molécule.
Ce chapitre construit ces orbitales à partir des \omterm{def:b2:fragment-orbitals:fragment}{orbitales de fragments}, explique
pourquoi l'eau est coudée et l'éthène plan, et pourquoi un carbocation préfère être entouré
de groupes alkyle.

\begin{recall}
\Cref{ch:b2:diatomic-mos}~: \omterm{def:b2:diatomic-mos:lcao}{CLOA}, recouvrement et symétrie, deux niveaux en interaction,
\omterm{def:b2:diatomic-mos:bonding}{orbitales liantes} et antiliantes. Le volume de première année~: géométries VSEPR, hybridation
comme description, mésomérie, carbocations et leur ordre de stabilité (l'interaction des
liaisons C--H avec l'orbitale vacante y était annoncée).
\end{recall}

\section{La méthode des fragments}

\begin{definition}[Fragments]\label{def:b2:fragment-orbitals:fragment}
Un \emph{fragment}\index{fragment} est un groupe d'atomes d'une molécule considéré
isolément (un atome, une paire d'hydrogènes, un groupe \ce{CH2}), et une \emph{orbitale de
fragment}\index{orbitale de fragment} l'une de ses orbitales. Les orbitales de la
molécule se construisent en faisant interagir les orbitales de deux fragments, deux à
deux, comme les orbitales atomiques d'une molécule diatomique.
\end{definition}

\begin{definition}[Combinaison adaptée à la symétrie]\label{def:b2:fragment-orbitals:symmetry-adapted}
Une \emph{combinaison adaptée à la symétrie}\index{combinaison adaptee a la symetrie@combinaison adaptée à la symétrie} d'orbitales
atomiques équivalentes est une combinaison que chaque opération de symétrie de la
molécule (réflexion dans un plan de symétrie, rotation autour d'un axe) soit
laisse inchangée, soit change en son opposée, ou, pour les ensembles dégénérés, transforme
en combinaisons du même ensemble.
\end{definition}

\begin{proposition}[Deux hydrogènes]\label{prop:b2:fragment-orbitals:h2-pair}
Pour deux hydrogènes équivalents d'orbitales 1s $h_1$ et $h_2$, les
\omterm{def:b2:fragment-orbitals:symmetry-adapted}{combinaisons adaptées à la symétrie} sont $h_1 + h_2$ et $h_1 - h_2$~: la première est
inchangée, la seconde change de signe, par toute opération qui échange les deux
atomes.
\end{proposition}

\begin{proof}
Une opération qui échange les atomes change $h_1$ en $h_2$ et $h_2$ en
$h_1$~: $h_1 + h_2$ est inchangée et $h_1 - h_2$ devient $h_2 - h_1 = -(h_1 - h_2)$.
\end{proof}

\begin{proposition}[Trois et quatre hydrogènes]\label{prop:b2:fragment-orbitals:h4}
Pour quatre hydrogènes aux sommets d'un tétraèdre, les combinaisons sont une
combinaison totalement symétrique $h_1 + h_2 + h_3 + h_4$ et un ensemble de trois combinaisons
dégénérées possédant chacune un plan nodal, en forme d'orbitales p. Pour trois
hydrogènes formant un triangle (comme dans l'ammoniac), une combinaison symétrique et une
paire dégénérée.
\end{proposition}

\admitted

Les formes des combinaisons dégénérées, et leurs noms (a, e, t), viennent de la théorie
des groupes, développée dans le volume de troisième année~; ici, ces noms ne sont que des étiquettes.

\begin{proposition}[Seul le semblable interagit avec le semblable]\label{prop:b2:fragment-orbitals:interaction-rules}
Une \omterm{def:b2:fragment-orbitals:fragment}{orbitale de fragment} n'interagit qu'avec les orbitales de l'autre \omterm{def:b2:fragment-orbitals:fragment}{fragment} qui
se comportent de la même façon sous chaque opération de symétrie de la molécule~; la
force de l'interaction croît avec le recouvrement et décroît avec l'écart
d'énergie.
\end{proposition}

\begin{proof}
Si deux orbitales se comportent différemment par une réflexion, l'une est symétrique et
l'autre antisymétrique par rapport à ce plan, et leur recouvrement et leur \omterm{def:b2:diatomic-mos:integrals}{intégrale de
résonance} s'annulent (\cref{prop:b2:diatomic-mos:symmetry}). La dépendance vis-à-vis du
recouvrement et de l'écart est celle du \cref{thm:b2:diatomic-mos:heteronuclear}.
\end{proof}

\begin{method}[Un diagramme de fragments]\label{met:b2:fragment-orbitals:fragment-diagram}
\begin{enumerate}
  \item Découper la molécule en deux \omterm{def:b2:fragment-orbitals:fragment}{fragments}, en général un atome central et l'ensemble
        des atomes qui l'entourent.
  \item Construire les \omterm{def:b2:fragment-orbitals:symmetry-adapted}{combinaisons adaptées à la symétrie} des orbitales périphériques.
  \item Classer les orbitales des deux \omterm{def:b2:fragment-orbitals:fragment}{fragments} selon leur comportement par rapport aux plans
        de symétrie et aux axes de la molécule.
  \item Faire interagir chaque orbitale avec celles de même comportement, d'autant plus fortement
        qu'elles sont proches en énergie~; celles qui restent sont non liantes.
  \item Remplir avec les électrons de valence et lire~: plus haut niveau occupé,
        caractère liant, préférences géométriques.
\end{enumerate}
\end{method}

\section{L'eau}

Plaçons l'eau dans le plan $yz$, l'axe $z$ étant la bissectrice de l'angle H--O--H. Deux
plans de symétrie contiennent l'axe $z$~: le plan moléculaire $yz$ et le plan $xz$
qui lui est perpendiculaire. Par rapport au plan $xz$, qui échange les deux
hydrogènes, $h_1 + h_2$ est symétrique et $h_1 - h_2$ antisymétrique. Sur l'oxygène, 2s et
$2\mathrm p_z$ sont symétriques par rapport aux deux plans~; $2\mathrm p_y$ (dans le
plan moléculaire, perpendiculaire à $z$) est antisymétrique par rapport à $xz$~;
$2\mathrm p_x$ est antisymétrique par rapport au plan moléculaire, ce qu'aucune
combinaison des hydrogènes n'est.

Ainsi, $h_1 + h_2$ se mélange avec 2s et $2\mathrm p_z$, $h_1 - h_2$ avec $2\mathrm p_y$, et
$2\mathrm p_x$ reste non liante. Il en résulte six orbitales, dont quatre sont remplies par les huit
électrons de valence.

\begin{omfigure}
\begin{tikzpicture}[font=\small]
  % O fragment (left)
  \pic (os) at (0,-2.4) {omlevel={ud}{0.7}}; \node[left] at (-0.45,-2.4) {2s};
  \pic (op1) at (-0.45,0) {omlevel={ud}{0.4}}; \pic (op2) at (0,0) {omlevel={u}{0.4}}; \pic (op3) at (0.45,0) {omlevel={u}{0.4}};
  \node[left] at (-0.7,0) {2p};
  \node at (0,-3.2) {O};
  % H2 fragment (right)
  \pic (hp) at (6,0.9) {omlevel={u}{0.7}}; \node[right] at (6.4,0.9) {$h_1 + h_2$};
  \pic (hm) at (6,1.8) {omlevel={u}{0.7}}; \node[right] at (6.4,1.8) {$h_1 - h_2$};
  \node[anchor=north west] at (6.4,0.6) {fragment \ce{H2}};
  % MOs
  \pic (m1) at (3,-2.9) {omlevel={ud}{0.7}};  \node[omsymlabel, below=0.17cm, inner sep=0.5pt] at (m1-c) {$2a_1$};
  \pic (m2) at (3,-1.55) {omlevel={ud}{0.7}};  \node[omsymlabel, below=0.17cm, inner sep=0.5pt] at (m2-c) {$1b_2$};
  \pic (m3) at (3,-0.78) {omlevel={ud}{0.7}};  \node[omsymlabel, below=0.17cm, inner sep=0.5pt] at (m3-c) {$3a_1$};
  \pic (m4) at (3,0) {omlevel={ud}{0.7}};  \node[omsymlabel, below=0.17cm, inner sep=0.5pt] at (m4-c) {$1b_1$};
  \pic (m5) at (3,2.4) {omlevel={empty}{0.7}}; \node[omsymlabel, below=0.17cm, inner sep=0.5pt] at (m5-c) {$4a_1$};
  \pic (m6) at (3,3.1) {omlevel={empty}{0.7}}; \node[omsymlabel, below=0.17cm, inner sep=0.5pt] at (m6-c) {$2b_2$};
  \draw[omcorr, densely dashed] (os-r) -- (m1-l) (op3-r) -- (m2-l) (op3-r) -- (m3-l) (op3-r) -- (m4-l) (op3-r) -- (m5-l) (op3-r) -- (m6-l);
  \draw[omcorr, densely dashed] (hp-l) -- (m1-r) (hm-l) -- (m2-r) (hp-l) -- (m3-r) (hp-l) -- (m5-r) (hm-l) -- (m6-r);
\end{tikzpicture}

{\small Diagramme de \omterm{def:b2:fragment-orbitals:fragment}{fragments} de l'eau, schéma. La combinaison symétrique des
hydrogènes se mélange avec les orbitales 2s et $2\mathrm p_z$ de l'oxygène (orbitales $a_1$),
l'antisymétrique avec $2\mathrm p_y$ ($b_2$)~; $2\mathrm p_x$, perpendiculaire au
plan moléculaire, reste non liante ($1b_1$, au niveau de la 2p de l'oxygène). Quatre niveaux de valence
remplis, quatre bandes photoélectroniques~; le plus haut, $1b_1$, est ionisé à
\qty{12.62}{eV}. Les étiquettes
$a_1$, $b_1$, $b_2$ sont des noms (issus de la théorie des groupes, vue dans le volume de troisième année).}
% ledger: ie:H2O
\end{omfigure}

\begin{proposition}[Pourquoi l'eau est coudée]\label{prop:b2:fragment-orbitals:bent-water}
Dans une molécule H--O--H linéaire, deux orbitales p de l'oxygène perpendiculaires à l'axe sont
non liantes. La courbure permet à l'une d'elles, celle qui est dans le plan moléculaire, de recouvrir
$h_1 + h_2$ et de se mélanger avec l'orbitale 2s~: ce niveau rempli descend, de plus que
les autres niveaux remplis ne montent. Avec huit électrons de valence, la molécule coudée est
la plus stable.
\end{proposition}

\begin{proof}[Argument]
Dans la molécule linéaire ($z$ selon H--O--H), $h_1 + h_2$ a un recouvrement nul avec
$2\mathrm p_x$ et $2\mathrm p_y$ (chacune est antisymétrique par rapport à un plan
contenant l'axe, la combinaison symétrique)~; ce sont deux doublets non liants
dégénérés. Quand on courbe la molécule dans le plan $yz$, $2\mathrm p_y$ pointe en partie vers
les hydrogènes~: son recouvrement avec $h_1 + h_2$ croît à partir de zéro, et le niveau rempli
qu'elle porte est stabilisé (le niveau $3a_1$ de la molécule coudée). Le niveau liant
$1b_2$ perd un peu de recouvrement et monte légèrement. Le bilan est un abaissement,
maximal lorsque le niveau qui descend est rempli, comme c'est le cas avec huit électrons. Une
molécule à quatre électrons de valence seulement, comme \ce{BeH2}, a ce niveau
vide et est linéaire.
\end{proof}

\section{Méthane, ammoniac et spectres photoélectroniques}

Dans le méthane, les quatre orbitales des hydrogènes donnent la combinaison symétrique, qui
correspond à l'orbitale 2s du carbone, et un ensemble dégénéré de trois, qui correspond à
ses trois orbitales 2p. Quatre \omterm{def:b2:diatomic-mos:bonding}{orbitales liantes} sont remplies~: une basse ($a_1$) et trois
dégénérées plus haut en énergie ($t_2$). L'ammoniac, avec trois hydrogènes, garde sur l'azote une
orbitale remplie orientée à l'opposé des hydrogènes, son doublet non liant, comme
plus haut niveau occupé.

\begin{omfigure}
\begin{tikzpicture}[font=\small]
  \pic (cs) at (0,-1.6) {omlevel={ud}{0.7}}; \node[left] at (-0.45,-1.6) {2s};
  \pic (cp1) at (-0.45,0) {omlevel={u}{0.4}}; \pic (cp2) at (0,0) {omlevel={u}{0.4}}; \pic (cp3) at (0.45,0) {omlevel={empty}{0.4}};
  \node[left] at (-0.7,0) {2p};
  \node at (0,-2.5) {C};
  \pic (ha) at (6,-0.8) {omlevel={ud}{0.7}}; \node[right] at (6.4,-0.8) {$a_1$};
  \pic (ht1) at (5.55,0.4) {omlevel={u}{0.4}}; \pic (ht2) at (6,0.4) {omlevel={u}{0.4}}; \pic (ht3) at (6.45,0.4) {omlevel={empty}{0.4}};
  \node[right] at (6.7,0.4) {$t_2$};
  \node at (6,-1.6) {fragment \ce{H4}};
  \pic (m1) at (3,-2.4) {omlevel={ud}{0.7}}; \node[omsymlabel, below=0.17cm, inner sep=0.5pt] at (m1-c) {$1a_1$};
  \pic (m2a) at (2.55,-0.9) {omlevel={ud}{0.4}}; \pic (m2b) at (3,-0.9) {omlevel={ud}{0.4}}; \pic (m2c) at (3.45,-0.9) {omlevel={ud}{0.4}};
  \node[omsymlabel, below=0.17cm, inner sep=0.5pt] at (m2b-c) {$1t_2$};
  \pic (m3) at (3,1.8) {omlevel={empty}{0.7}}; \node[omsymlabel, below=0.17cm, inner sep=0.5pt] at (m3-c) {$2a_1$};
  \pic (m4a) at (2.55,2.5) {omlevel={empty}{0.4}}; \pic (m4b) at (3,2.5) {omlevel={empty}{0.4}}; \pic (m4c) at (3.45,2.5) {omlevel={empty}{0.4}};
  \node[omsymlabel, below=0.17cm, inner sep=0.5pt] at (m4b-c) {$2t_2$};
  \draw[omcorr, densely dashed] (cs-r) -- (m1-l) (cp3-r) -- (m2a-l) (cs-r) -- (m3-l) (cp3-r) -- (m4a-l);
  \draw[omcorr, densely dashed] (ha-l) -- (m1-r) (ht1-l) -- (m2c-r) (ha-l) -- (m3-r) (ht1-l) -- (m4c-r);
\end{tikzpicture}

{\small Diagramme de \omterm{def:b2:fragment-orbitals:fragment}{fragments} du méthane, schéma. La combinaison symétrique des
quatre hydrogènes correspond à l'orbitale 2s du carbone, les trois combinaisons dégénérées à ses
orbitales 2p~: deux niveaux remplis, donc deux bandes photoélectroniques, la plus haute
(ionisée la première, à \qty{12.61}{eV}) contenant six électrons.}
% ledger: ie:CH4
\end{omfigure}

\begin{definition}[Spectroscopie photoélectronique]\label{def:b2:fragment-orbitals:photoelectron}
La \emph{spectroscopie photoélectronique}\index{spectroscopie photoelectronique@spectroscopie photoélectronique} mesure les
énergies cinétiques des électrons arrachés à des molécules par un rayonnement
monochromatique d'énergie connue~; chaque bande du spectre correspond à l'arrachement
d'un électron d'un niveau occupé, et sa position donne l'énergie d'ionisation
de ce niveau.
\end{definition}

\begin{proposition}[Approximation de Koopmans]\label{prop:b2:fragment-orbitals:koopmans}
L'énergie nécessaire pour arracher un électron d'une orbitale occupée est
approximativement l'opposé de l'énergie de cette orbitale, $I \approx -\varepsilon$.
\end{proposition}

\admitted

L'approximation néglige la relaxation des autres électrons après
l'ionisation~; elle est justifiée dans le volume de troisième année. Elle transforme un spectre
photoélectronique en diagramme d'orbitales~: l'eau présente quatre bandes de valence, le méthane deux,
en accord avec les diagrammes de \omterm{def:b2:fragment-orbitals:fragment}{fragments} --- et contre l'image de deux doublets non liants
équivalents et de deux liaisons équivalentes dans l'eau, qui décrit la même densité électronique
totale d'une autre manière mais ne donne pas les énergies des
électrons arrachés.

\section{L'éthène}

On peut construire l'éthène à partir de deux \omterm{def:b2:fragment-orbitals:fragment}{fragments} \ce{CH2}. Chaque \ce{CH2} possède, outre les
orbitales de ses deux liaisons C--H, une orbitale orientée à l'opposé des hydrogènes dans le
plan (qui forme la liaison $\sigma$ C--C avec son partenaire) et une orbitale p
perpendiculaire au plan. Les deux orbitales p se combinent en une orbitale $\pi$ pleine et une
orbitale $\pi^*$ vide.

\begin{omfigure}
\begin{tikzpicture}[font=\small]
  \begin{scope}
    \pic at (0,0) {omp={0.85}{90}}; \pic at (1.2,0) {omp={0.85}{90}};
    \draw[thick] (0,0) -- (1.2,0);
    \draw[thick] (0,0) -- (-0.5,-0.45) (0,0) -- (-0.55,0.35) (1.2,0) -- (1.7,-0.45) (1.2,0) -- (1.75,0.35);
    \node at (0.6,-1.1) {$\pi$ (pleine)};
  \end{scope}
  \begin{scope}[xshift=3.6cm]
    \pic at (0,0) {omp={0.85}{90}}; \pic at (1.2,0) {omp={-0.85}{90}};
    \draw[thick] (0,0) -- (1.2,0);
    \draw[thick] (0,0) -- (-0.5,-0.45) (0,0) -- (-0.55,0.35) (1.2,0) -- (1.7,-0.45) (1.2,0) -- (1.75,0.35);
    \node at (0.6,-1.1) {$\pi^*$ (vide)};
  \end{scope}
  \begin{scope}[xshift=7.4cm]
    \pic at (0,0) {omp={0.85}{90}}; \pic at (1.2,0) {omp={0.85}{0}};
    \draw[thick] (0,0) -- (1.2,0);
    \node at (0.6,-1.1) {tourné de $90^\circ$};
    \node[font=\scriptsize, align=center] at (0.6,1.3) {recouvrement nul};
  \end{scope}
\end{tikzpicture}

{\small Le système $\pi$ de l'éthène, schéma (la molécule vue par la tranche, les
hydrogènes dans le plan). Les deux orbitales p perpendiculaires au plan se combinent
en phase ($\pi$, pleine) et en opposition de phase ($\pi^*$, vide). Tourner un \ce{CH2}
de $90^\circ$ rend les deux orbitales p perpendiculaires~: leur recouvrement, et la liaison $\pi$,
s'annulent.}
\end{omfigure}

La liaison $\pi$ maintient l'éthène plan~: tourner une extrémité écarte les orbitales p l'une de
l'autre, leur recouvrement décroît comme le cosinus de l'angle de torsion, et à $90^\circ$
la liaison $\pi$ a disparu. Il faut fournir l'énergie d'une liaison $\pi$ pour tordre la
molécule~: c'est pourquoi les isomères cis et trans des alcènes ne s'interconvertissent pas à
température ambiante, alors que la rotation autour d'une liaison simple est libre.

\section{L'hyperconjugaison}

\begin{definition}[Hyperconjugaison]\label{def:b2:fragment-orbitals:hyperconjugation}
L'\emph{hyperconjugaison}\index{hyperconjugaison} est l'interaction d'une orbitale
$\sigma$ pleine d'une liaison C--H ou C--C avec une orbitale p ou $\pi^*$ voisine, vacante ou
partiellement remplie, qui lui est parallèle.
\end{definition}

\begin{proposition}[Les groupes alkyle stabilisent les carbocations]\label{prop:b2:fragment-orbitals:hyperconjugation}
Un carbocation est stabilisé par chaque liaison C--H ou C--C portée par un carbone voisin
qui peut s'aligner avec son orbitale p vacante~: plus il y a de telles liaisons, plus le
cation est stable, d'où l'ordre tertiaire > secondaire > primaire > méthyle.
\end{proposition}

\begin{proof}[Argument]
L'orbitale p vacante du carbone cationique et une orbitale $\sigma_{\text{C--H}}$ pleine
voisine forment un système à deux niveaux avec un grand écart $\Delta$ et un faible couplage
$\beta$ (non nul seulement si la liaison est à peu près parallèle à l'orbitale p). D'après le
\cref{thm:b2:diatomic-mos:heteronuclear}, le niveau $\sigma$ plein est abaissé d'environ
$\beta^2/\Delta$, et les deux électrons gagnent le double~; le niveau vacant monte, sans
coût. Chaque liaison convenablement alignée ajoute sa propre stabilisation~; un groupe méthyle
a toujours une liaison C--H presque alignée, quelle que soit sa rotation.
\end{proof}

\begin{omfigure}
\begin{tikzpicture}[font=\small]
  \pic at (0,0) {omp={0.9}{90}};
  \node[circle, fill=black, inner sep=1.2pt] at (0,0) {};
  \node[font=\scriptsize, anchor=north west] at (0.08,-0.05) {C$^+$};
  \draw[thick] (0,0) -- (1.4,0);
  \node[circle, fill=black, inner sep=1.2pt] at (1.4,0) {};
  \draw[thick] (1.4,0) -- (1.75,0.85);
  \pic at (1.58,0.42) {omlobe={0.75}{68}};
  \pic at (1.58,0.42) {omlobe={0.45}{248}};
  \node[font=\scriptsize] at (1.95,1.0) {H};
  \draw[thick] (1.4,0) -- (2.0,-0.45) (1.4,0) -- (1.1,-0.6);
  \draw[<->, omProp, thick] (0.35,0.8) .. controls (0.8,1.2) and (1.2,1.2) .. (1.5,0.95);
  \node[font=\scriptsize, align=left, anchor=west] at (2.4,0.4) {lobe $\sigma_{\text{C--H}}$ plein parallèle\\ à l'orbitale p vacante};
  % level scheme
  \begin{scope}[xshift=7.2cm, yshift=-0.4cm]
    \pic (p) at (0,1.4) {omlevel={empty}{0.7}}; \node[left] at (-0.4,1.4) {p vacante};
    \pic (s) at (2.2,-0.8) {omlevel={ud}{0.7}}; \node[right] at (2.6,-0.8) {$\sigma_{\text{C--H}}$};
    \pic (lo) at (1.1,-1.2) {omlevel={ud}{0.7}};
    \pic (hi) at (1.1,1.7) {omlevel={empty}{0.7}};
    \draw[omcorr, densely dashed] (p-r) -- (lo-l) (p-r) -- (hi-l) (s-l) -- (lo-r) (s-l) -- (hi-r);
    \node[font=\scriptsize, anchor=north] at (1.1,-1.35) {stabilisé de $\approx \beta^2/\Delta$};
  \end{scope}
\end{tikzpicture}

{\small \omterm{def:b2:fragment-orbitals:hyperconjugation}{Hyperconjugaison} dans un carbocation, schéma. Une \omterm{def:b2:diatomic-mos:bonding}{orbitale liante} C--H pleine
portée par le carbone voisin, parallèle à l'orbitale p vacante, lui cède une partie de sa
densité électronique~; le niveau plein est stabilisé comme dans toute interaction de
deux niveaux d'énergies différentes.}
\end{omfigure}

\section{Exercices}

\begin{exercise}[$\star$]\label{exo:b2:fragment-orbitals:1}
Écrire les deux \omterm{def:b2:fragment-orbitals:symmetry-adapted}{combinaisons adaptées à la symétrie} des orbitales 1s des hydrogènes de
l'eau et dire laquelle change de signe par réflexion dans le plan qui contient la bissectrice de l'angle H--O--H
et qui est perpendiculaire à la molécule.
\end{exercise}

\begin{exercise}[$\star$]\label{exo:b2:fragment-orbitals:2}
Combien d'\omterm{def:b2:diatomic-mos:lcao}{orbitales moléculaires} de valence l'eau possède-t-elle, combien d'électrons de valence,
et combien d'orbitales de valence remplies~?
\end{exercise}

\begin{exercise}[$\star$]\label{exo:b2:fragment-orbitals:3}
Pourquoi le spectre photoélectronique du méthane présente-t-il deux bandes de valence, l'une
trois fois plus intense que l'autre~?
\end{exercise}

\begin{exercise}[$\star$]\label{exo:b2:fragment-orbitals:4}
Pourquoi les six atomes de l'éthène sont-ils dans un même plan~?
\end{exercise}

\begin{exercise}[$\star\star$]\label{exo:b2:fragment-orbitals:5}
L'énergie de première ionisation de l'ammoniac vaut \qty{10.07}{eV}, celle du méthane
\qty{12.61}{eV}. Quelle orbitale perd l'électron dans chaque cas, et pourquoi l'ammoniac est-il le
plus facile à ioniser~?
\end{exercise}

\begin{exercise}[$\star\star$]\label{exo:b2:fragment-orbitals:6}
Comparer les niveaux occupés de l'eau linéaire et de l'eau coudée, et expliquer pourquoi une molécule
comme \ce{BeH2}, à quatre électrons de valence, est linéaire.
\end{exercise}

\begin{exercise}[$\star\star$]\label{exo:b2:fragment-orbitals:7}
Expliquer, à l'aide du résultat à deux niveaux, l'ordre de stabilité des carbocations
\ce{CH3+}, \ce{CH3CH2+}, \ce{(CH3)2CH+}, \ce{(CH3)3C+}.
\end{exercise}

\begin{exercise}[$\star\star$]\label{exo:b2:fragment-orbitals:8}
Le recouvrement de deux orbitales p tournées d'un angle $\theta$ autour de la liaison est
proportionnel à $\cos\theta$. L'\omterm{def:b2:diatomic-mos:integrals}{intégrale de résonance} étant proportionnelle au
recouvrement, comment la stabilisation $\pi$ dépend-elle de $\theta$~? Pour quel angle tombe-t-elle
à la moitié~?
\end{exercise}

\begin{exercise}[$\star\star$]\label{exo:b2:fragment-orbitals:9}
Le borane \ce{BH3} est plan, l'ammoniac pyramidal. À l'aide des \omterm{def:b2:fragment-orbitals:fragment}{orbitales de fragments} d'un
atome central et de trois hydrogènes, expliquer la différence par les électrons dont chacun dispose.
\end{exercise}

\begin{exercise}[$\star\star\star$]\label{exo:b2:fragment-orbitals:10}
L'ion \ce{H3+} est un triangle équilatéral. Avec $\alpha$ et $\beta$ pour les orbitales 1s
et en négligeant le recouvrement, trouver ses trois \omterm{def:b2:atomic-orbitals:orbital-energy}{énergies orbitalaires} (la combinaison
symétrique a pour énergie $\alpha + 2\beta$) et expliquer pourquoi deux électrons suffisent à
le lier.
\end{exercise}

\begin{exercise}[$\star\star\star$]\label{exo:b2:fragment-orbitals:11}
Le système allyle \ce{CH2=CH-CH2} possède trois orbitales p parallèles. Deviner la forme de
ses trois orbitales $\pi$ (signes sur chaque carbone) d'après le nombre de nœuds, et dire laquelle
est le plus haut niveau occupé de l'anion allyle.
\end{exercise}

\begin{exercise}[$\star\star\star$]\label{exo:b2:fragment-orbitals:12}
Construire le système $\pi$ du dioxyde de carbone à partir des orbitales 2p perpendiculaires à
l'axe O--C--O~: combien d'orbitales $\pi$, lesquelles sont liantes, non liantes,
antiliantes, et combien d'électrons contiennent-elles~?
\end{exercise}

\section{Problème~: la forme de l'eau}

\begin{problem}[{Problème du week-end --- les \omterm{def:b2:fragment-orbitals:fragment}{fragments} oxygène et \ce{H2}, les orbitales
d'une molécule H--O--H linéaire, les niveaux qui se déplacent quand elle se courbe, et le spectre
photoélectronique avec l'approximation de Koopmans}]
\label{pb:b2:fragment-orbitals:1}
Données~: géométrie de l'eau, $r(\ce{O-H}) = \qty{95.8}{pm}$, angle
$104.48^\circ$~; énergies de première ionisation~: \ce{H2O} \qty{12.62}{eV}, \ce{O}
\qty{13.62}{eV}. L'eau est dans le plan $yz$, $z$ selon la bissectrice.
% ledger: geo:H2O.rOH, geo:H2O.aHOH, ie:H2O, ie:O

\medskip
\noindent\textbf{Partie I --- Les \omterm{def:b2:fragment-orbitals:fragment}{fragments}.}
\begin{enumerate}
  \item Combien d'électrons de valence l'eau possède-t-elle, et de quels atomes proviennent-ils~?
  \item Écrire les deux \omterm{def:b2:fragment-orbitals:symmetry-adapted}{combinaisons adaptées à la symétrie} des orbitales 1s des hydrogènes.
  \item Quelles orbitales de l'oxygène sont symétriques par rapport aux deux plans de symétrie~?
  \item Quelle orbitale de l'oxygène est antisymétrique par rapport au seul plan $xz$~?
  \item Quelle orbitale de l'oxygène n'a aucun partenaire parmi les combinaisons des hydrogènes~?
  \item Combien d'\omterm{def:b2:diatomic-mos:lcao}{orbitales moléculaires} en résulte-t-il, et combien sont remplies~?
\end{enumerate}

\medskip
\noindent\textbf{Partie II --- Une molécule H--O--H linéaire.}
\begin{enumerate}[resume]
  \item Dans une molécule linéaire selon $z$, quelle combinaison interagit avec la
        $2\mathrm p_z$ de l'oxygène~?
  \item Laquelle interagit avec la 2s de l'oxygène~?
  \item Quelles orbitales de l'oxygène restent non liantes, et avec quelle dégénérescence~?
  \item Remplir les niveaux avec les huit électrons de valence.
  \item L'eau linéaire serait-elle \omterm{def:b2:diatomic-mos:magnetism}{paramagnétique}~?
  \item Où se situerait son plus haut niveau occupé, par rapport à une orbitale 2p de
        l'oxygène~?
\end{enumerate}

\medskip
\noindent\textbf{Partie III --- La courbure.}
\begin{enumerate}[resume]
  \item La molécule linéaire couchée selon $z$ est courbée dans le plan $yz$. Quelle
        \omterm{def:b2:diatomic-mos:bonding}{orbitale non liante} commence à recouvrir $h_1 + h_2$~?
  \item Que devient le niveau rempli qu'elle porte~?
  \item Quel niveau liant monte légèrement, et pourquoi~?
  \item Conclure~: pourquoi l'eau est-elle coudée~?
  \item Comparer l'angle mesuré à $90^\circ$ (liaison p pure) et à $109.5^\circ$.
  \item Quelle \omterm{def:b2:diatomic-mos:bonding}{orbitale non liante} la courbure laisse-t-elle intacte~?
\end{enumerate}

\medskip
\noindent\textbf{Partie IV --- Le spectre photoélectronique.}
\begin{enumerate}[resume]
  \item Combien de bandes de valence le spectre présente-t-il~?
  \item D'après l'approximation de Koopmans, de quelle orbitale provient la première bande~?
  \item Pourquoi cette bande est-elle étroite, alors que les bandes des \omterm{def:b2:diatomic-mos:bonding}{orbitales liantes} sont larges~?
  \item Comparer l'énergie de première ionisation de l'eau à celle de l'atome
        d'oxygène, et expliquer la différence.
  \item Que deviennent les «~deux doublets non liants équivalents~» du schéma de Lewis~?
  \item Donner le nombre de bandes d'ionisation de valence de l'eau et l'énergie de
        la plus basse.
\end{enumerate}
\end{problem}
